%% main.tex, Paper 2: Spectral Surrogates for Rough Stochastic Volatility
%%
%% Working title: Distributional Learning of Rough Volatility via Spectral
%%        Surrogates: a Cramér–Bochner PINN for fBm-Driven SDEs
%%
%% Target: SIAM Journal on Financial Mathematics / Quantitative Finance
%%
%% Build: pdflatex main → bibtex main → pdflatex main ×2
%% ─────────────────────────────────────────────────────────────────────────────

\documentclass[11pt,a4paper]{article}

\usepackage{amsmath, amssymb, amsthm}
\usepackage{mathtools}
\usepackage[margin=2.5cm]{geometry}
\usepackage{hyperref}
\usepackage{cleveref}
\usepackage{booktabs}
\usepackage{enumitem}
\usepackage{graphicx}
\graphicspath{{../02_rough_vol/figures/}}
\usepackage{xcolor}
\usepackage{microtype}
\usepackage{natbib}

% ── Theorem environments ──────────────────────────────────────────────────────
\newtheorem{theorem}{Theorem}[section]
\newtheorem{lemma}[theorem]{Lemma}
\newtheorem{corollary}[theorem]{Corollary}
\newtheorem{proposition}[theorem]{Proposition}
\newtheorem{remark}[theorem]{Remark}
\newtheorem{assumption}{Assumption}
\renewcommand{\theassumption}{A\arabic{assumption}}

% ── Notation ─────────────────────────────────────────────────────────────────
\newcommand{\R}{\mathbb{R}}
\newcommand{\E}{\mathbb{E}}
\newcommand{\Prob}{\mathbb{P}}
\newcommand{\Law}{\mathcal{L}}
\newcommand{\W}{W_2}        % Wasserstein-2
\newcommand{\Hcal}{\mathcal{H}}
\newcommand{\Lcal}{\mathcal{L}}
\newcommand{\norm}[1]{\|#1\|}
\newcommand{\HS}{\mathrm{HS}}
\newcommand{\dt}{\,\mathrm{d}t}
\newcommand{\dw}{\,\mathrm{d}W}
\newcommand{\TODO}[1]{\textcolor{red}{[\textbf{TODO:} #1]}}

% ─────────────────────────────────────────────────────────────────────────────
\begin{document}

\title{Distributional Learning of Rough Volatility via Spectral Surrogates:\\
    a Cramér PINN for Fractional Brownian Motion--Driven SDEs}

\author{Hubert Lipski}
\date{\today}
\maketitle

% ─────────────────────────────────────────────────────────────────────────────
\begin{abstract}
Rough stochastic volatility models, in which the log-variance is driven by a
fractional Brownian motion with Hurst exponent $H \approx 0.1$, reproduce
key empirical features of equity implied volatility but require expensive
quasi-Monte Carlo simulation for option pricing. We introduce the
\emph{Cramér PINN}, a neural-network surrogate trained to match the
power-spectral density (PSD) of fractional Gaussian noise (fGn), the
stationary increment process of fBm. The surrogate requires no
numerical solver at inference time and generates variance paths by a
single forward pass. We prove a $W_2$ convergence theorem showing that
matching the exact fGn PSD forces the surrogate's law to converge to the
reference process, with an explicit Hurst-exponent consistency corollary:
if the training loss is $\varepsilon$, then $|\hat{H} - H| \leq C\sqrt{\varepsilon}$.
On the rough Bergomi model ($H=0.1$, $\eta=1.9$, $\xi_0=0.04$), a linear
spectral filter surrogate trained on the exact fGn PSD achieves an
implied-volatility surface RMSE of \textbf{1.33~vol~points} vs.\ the MC
reference, compared to \textbf{2.27~vp} for the asymptotic
$C|f|^{1-2H}$ approximation. A nonlinear MLP surrogate trained with
the same spectral loss reaches \textbf{1.21~vp}, confirming that the loss
is architecture-agnostic. Under rough Heston with $\rho=-0.9$, we show that any surrogate matching
only the marginal distribution of $V$ carries an \textbf{irreducible ATM
delta bias of $-0.27$} ($35\%$ relative). We further show that circular
FFT surrogates fail catastrophically at $\rho\neq 0$ (204~vp, delta $>1$)
due to non-causal wrap-around, and that a \textbf{CausalSpectralFilter}
(lower-triangular Toeplitz convolution, same PSD loss) resolves this,
achieving \textbf{3.26~vp} vs.\ the BLP reference and reducing ATM delta
bias to $-0.07$. Training the causal filter on the reference's own
RL-fBm PSD eliminates the residual entirely: \textbf{2.57~vp} and ATM
delta bias $\lesssim 0.01$ (Monte-Carlo-noise level). On rough Heston,
the same spectral causal architecture cuts the cross-model bias from
$+0.49$ to $-0.15$; a Volterra-aware autoregressive surrogate with
$\sqrt{V_j}$ diffusion closes the residual to Monte-Carlo noise
(\textbf{0.80~vp}, bias $\leq 0.005$). The same architecture, trained
end-to-end from a misspecified init, calibrates to a live SPY option
chain across three expirations ($32$, $88$, $179$ days; $26$ strikes
total) at \textbf{$3.97$~vp overall IV RMSE} in $3$~minutes on a
single GPU, with shared parameters across maturities.
\end{abstract}

% ═══════════════════════════════════════════
\section{Introduction}
\label{sec:intro}

This paper is the second part of a study on spectral surrogates for
rough volatility. Paper~2a \citep{TODO-paper2a} introduced the
Cramér PINN, a neural-network surrogate for the rough Bergomi (rBergomi)
model trained to match the exact PSD of fractional Gaussian noise.
It established a $W_2$ convergence theorem, validated SpectralFilter and
CausalSpectralFilter at $\rho=0$ and $\rho=-0.9$, and demonstrated
that matching the training PSD to the reference simulator eliminates
residual delta bias to Monte Carlo noise.

In this second part, we address the cross-model problem: the CausalSpectralFilter
resolves $\rho\neq 0$ within rBergomi, but when the reference process is
rough Heston (with mean-reverting $V$ dynamics), a residual delta bias of
$-0.15$ remains. We introduce the VolterraSurrogate, an autoregressive
architecture with $\sqrt{V_j}$ diffusion, and show that it closes the
cross-model gap to Monte Carlo noise ($0.80$~vp, delta bias $\leq 0.005$).
End-to-end training from misspecified initialisation, Markovian kernel
approximation ($7\times$ speedup), and live SPY calibration ($3.97$~vp,
$3$~minutes on a single GPU) confirm the method's practical viability.

\tableofcontents
\newpage
\subsection{Leverage effect and Greek bias ($\rho = -0.9$)}
\label{sec:greeks}

Real equity and index markets exhibit a strong negative correlation
$\rho \approx -0.9$ between price and volatility (the \emph{leverage
effect}). Incorporating $\rho \neq 0$ into the rBergomi model requires
the hybrid scheme of \citet{bennedsen2017hybrid}, which explicitly
separates the standard-BM innovations driving $W^H$ from the fBm
increments; naively correlating the stock with fGn increments introduces
a spurious predictable drift. We therefore conduct the $\rho \neq 0$
experiment on the \emph{rough Heston} model
\citep{ElEuch2019,el2019perfect}, which natively exposes the standard-BM
innovations $Z_k$ through its Volterra representation and admits a clean
Cholesky decomposition $\mathrm{d}Z^S = \rho Z_k + \sqrt{1-\rho^2}
Z_k^\perp$.

\paragraph{Setup.}
Rough Heston parameters:
$H = 0.1$, $\kappa = 2.0$, $\bar{V} = 0.04$, $\nu = 1.9$,
$V_0 = 0.04$, $\rho = -0.9$.
Reference: 8\,000 paths from a forward-Euler Volterra discretisation
(O($N^2$) per path, vectorised).
Three experiments are run with fixed reference variance paths
$V_{\mathrm{rH}}$ and BM innovations $Z_{\mathrm{rH}}$:
\begin{itemize}[nosep]
 \item \textbf{Reference}: $V_{\mathrm{rH}}$, $Z_{\mathrm{rH}}$
    for $\rho$, correct joint law.
 \item \textbf{Surrogate A}: $V_{\mathrm{rH}}$, fresh independent
    $Z_{\mathrm{fresh}}$ for $\rho$, \emph{same} variance
    distribution, wrong BM. Isolates the pure correlation bias.
 \item \textbf{Surrogate B}: SpectralFilter variance $V_{\mathrm{surr}}$
    (trained on rBergomi fGn PSD, $\rho = -0.9$), wrong variance
    distribution \emph{and} wrong BM.
 \item \textbf{Surrogate C}: $V_{\mathrm{surr}}$, $\rho = 0$ --
    baseline; no leverage modelled.
\end{itemize}

\paragraph{IV surface.}
Table~\ref{tab:iv-rh} shows the reference smile.
The leverage effect creates a strong left skew: IV falls from 36.1\% at
$K/S_0 = 0.80$ to 8.9\% at $K/S_0 = 1.05$ for $T = 0.25$, with a
slight uptick at deep OTM calls ($K/S_0 \geq 1.10$) due to residual
upside tail mass (a standard rough-vol feature).
Surrogate~C ($\rho = 0$) achieves an overall RMSE of 7.8~vp against the
reference, competitive, because the spectral surrogate captures the
\emph{level} of IV through the matched marginal distribution of $V$.
Surrogate~B, which naively adds $\rho = -0.9$ using the surrogate's
internal white noise, achieves 188~vp: using the wrong BM for the
Cholesky correlation is strictly worse than omitting correlation
entirely.

\begin{table}[h]
\centering
\caption{Rough Heston reference IV surface ($H=0.1$, $\kappa=2$,
$\nu=1.9$, $\rho=-0.9$, 8k paths). Left skew is pronounced;
SpectralFilter~(C, $\rho=0$) achieves overall RMSE 7.8~vp by
matching the level of the smile despite ignoring skew.}
\label{tab:iv-rh}
\begin{tabular}{lrrrrrrrrr}
\toprule
$T \backslash K/S_0$ & 0.80 & 0.85 & 0.90 & 0.95 & 1.00
           & 1.05 & 1.10 & 1.15 & 1.20 \\
\midrule
\multicolumn{10}{c}{\textit{Reference (rough Heston, $\rho=-0.9$)}} \\
0.25 & 36.1 & 30.8 & 25.4 & 19.7 & 13.5 & 8.9 & 9.4 & 10.7 & 10.4 \\
0.50 & 29.1 & 25.3 & 21.5 & 17.4 & 13.0 & 8.9 & 8.1 & 8.7 & 9.4 \\
1.00 & 24.7 & 21.8 & 19.0 & 16.1 & 13.2 & 10.2 & 8.1 & 7.8 & 8.1 \\
\midrule
\multicolumn{10}{c}{\textit{Surrogate C (SpectralFilter, $\rho=0$), RMSE 7.8~vp}} \\
0.25 & 28.1 & 24.3 & 21.0 & 18.4 & 17.2 & 18.1 & 20.0 & 21.7 & 23.5 \\
0.50 & 22.3 & 20.4 & 18.6 & 17.2 & 16.7 & 17.2 & 18.4 & 19.9 & 21.3 \\
1.00 & 20.4 & 19.0 & 18.0 & 17.4 & 17.2 & 17.5 & 18.2 & 19.1 & 20.1 \\
\bottomrule
\end{tabular}
\end{table}

\paragraph{Delta bias.}
Table~\ref{tab:delta-bias} and \cref{fig:delta-bias} report the ATM
($K/S_0 = 1.00$) delta and the bias
$\delta_\mathrm{bias} = \hat\Delta - \Delta_\mathrm{ref}$ across
experiments. Deltas are computed by bump-and-reprice ($\varepsilon =
1\%$) with frozen variance paths and frozen $Z$ innovations (common
random numbers).

\begin{table}[h]
\centering
\caption{ATM delta and bias under rough Heston ($\rho=-0.9$).
Bias~A uses the \emph{same} variance paths as the reference but an
independent BM for correlation; it isolates the pure leverage-effect
error. Bias~B uses the surrogate variance \emph{and} an independent
BM; it compounds variance-distribution and correlation errors.
All biases in $\Delta$-units; reference $\Delta \approx 0.78$ for all $T$.}
\label{tab:delta-bias}
\begin{tabular}{lrrrr}
\toprule
 & $\Delta_{\mathrm{ref}}$ & $\Delta_A$ & Bias~A & Bias~B \\
\midrule
$T = 0.25$ & 0.781 & 0.511 & $-0.270$ & $+0.290$ \\
$T = 0.50$ & 0.775 & 0.516 & $-0.258$ & $+0.488$ \\
$T = 1.00$ & 0.770 & 0.525 & $-0.245$ & $+1.159$ \\
\bottomrule
\end{tabular}
\end{table}

Three observations stand out.
\emph{(i) Reference delta is high and flat} ($\approx 0.78$ across all
maturities), well above the Black-Scholes ATM value ($\approx 0.51$).
This is the leverage-amplified right tail: when $S_T > K$, it is because
$V$ happened to stay low (due to $\rho = -0.9$), and such paths tend to
have $S_T \gg K$; the call payoff is large, inflating the sensitivity
$\partial C / \partial S_0$.
\emph{(ii) Bias~A is $\approx -0.27$, stable across all maturities.}
Breaking the $S$--$V$ dependence (replacing $Z_{\mathrm{rH}}$ with an
independent draw) removes the leverage amplification entirely, collapsing
delta to the symmetric value $\approx 0.51$. The surrogate therefore
under-estimates delta by \emph{35\%} for an ATM call under realistic
leverage. This bias is irreducible from marginal-distribution matching
alone: even a perfect replica of the variance marginal cannot recover it.
\emph{(iii) Bias~B grows with $T$} ($+0.29$ at $T = 0.25$,
$+1.16$ at $T = 1$) because the surrogate's rBergomi $V$ lacks the
mean-reversion of rough Heston: at long horizons $V_\mathrm{surr}$
remains elevated while $V_\mathrm{rH}$ reverts to $\bar V = 0.04$,
inflating option prices and delta beyond physically meaningful levels.

\begin{figure}[t]
 \centering
 \includegraphics[width=\textwidth]{rheston_iv_comparison}
 \caption{%
  \textbf{Rough Heston IV smiles: reference vs.\ surrogates.}
  Each panel shows one maturity.
  Reference (black solid) exhibits a pronounced left skew.
  Surrogate~B (red dashed, $\rho=-0.9$, wrong BM) is far from
  the reference; Surrogate~C (blue dotted, $\rho=0$) captures the
  level of the smile but not its skew. Overall RMSE: 7.8~vp
  (Surrogate~C) vs.\ 188~vp (Surrogate~B).
 }
 \label{fig:rh-iv}
\end{figure}

\begin{figure}[t]
 \centering
 \includegraphics[width=\textwidth]{rheston_delta_bias}
 \caption{%
  \textbf{Delta bias per strike and maturity.}
  Blue bars: Bias~A (same variance, fresh BM), the pure leverage
  contribution to delta that any marginal-matching surrogate misses.
  Orange bars: Bias~B (surrogate variance + fresh BM), compounded
  with the variance-distribution error, which dominates at $T=1$.
  Bias~A is negative for $K \leq 1.00$ (surrogate under-estimates
  delta for ITM/ATM calls) and slightly positive for deep OTM, a
  signature of the leverage effect on the conditional right-tail
  distribution.
 }
 \label{fig:delta-bias}
\end{figure}

\begin{remark}[Architecture scope and correcting the correlation bias]
\label{rem:rho}
 Bias~A is not an artefact of the surrogate architecture; it is a
 structural limitation of any model that matches only the marginal law
 of $V$. We verified this by attempting to train a SpectralFilter
 directly on the empirical PSD of rough Heston log-variance increments:
 the resulting surrogate collapses to $E[V_\text{surr}] \approx 0$
 because the exponential-of-fBm reconstruction
 $V = \xi_0 e^{\eta W - \frac{1}{2}\eta^2 \operatorname{Var}[W_t]}$
 is valid only for rBergomi (log-normal $V$) and cannot represent the
 mean-reverting rough Heston dynamics. This confirms that the Cramér
 PINN, as formulated, is an rBergomi surrogate; a fully consistent rough
 Heston surrogate requires either a Volterra-aware architecture or
 direct simulation.

 Correcting the remaining delta bias fully requires training on joint
 $(V, S)$ paths or a post-hoc Malliavin score correction; both are left
 for future work. Section~\ref{sec:causal} shows that the circular-FFT
 architecture itself can be replaced by a causal convolution that at
 least encodes $\rho \neq 0$ correctly within rBergomi.
\end{remark}

% ═════════════════════════════════════════════════════════════════════════════
\subsection{CausalSpectralFilter: supporting $\rho \neq 0$}
\label{sec:causal}
% ═════════════════════════════════════════════════════════════════════════════

The experiments in \cref{sec:greeks} establish that the irreducible
$-0.27$ delta bias originates from missing the joint $(V,S)$ law. A
natural prior question is whether the spectral surrogate can at least
\emph{encode} $\rho \neq 0$ correctly within rBergomi, so that the
leverage channel is present even if the exact $S$--$V$ law is not matched.

\paragraph{Setup: $z$-Cholesky and the BLP reference.}
The \emph{$z$-Cholesky} approach reuses the filter's i.i.d.\ input noise
$z \sim \mathcal{N}(0,I)$ as the stock BM driver:
\begin{equation}
 d\!\log S_k
  = -\tfrac{1}{2}V_k\,dt
   + \sqrt{V_k\,dt}\,\bigl(\rho\,z_k + \sqrt{1-\rho^2}\,z_k^\perp\bigr).
 \label{eq:z-cholesky}
\end{equation}
The ground truth is an RL-discretised rBergomi path with $\rho = -0.9$,
simulated via the \emph{BLP hybrid scheme} \citep{bennedsen2017hybrid}
(same Volterra kernel as \cref{sec:greeks}, now applied to rBergomi).
The BLP reference produces a left-skewed IV surface (ATM $\approx 15\%$
at $T=0.25$, declining to $9\%$ at $K/S_0 = 1.10$) and ATM delta
$0.67$--$0.71$, consistent with the leverage effect.

\paragraph{Circular filter: catastrophic failure.}
Applying \eqref{eq:z-cholesky} to the circular SpectralFilter gives an
IV RMSE of $204$~vp and ATM deltas rising from $1.07$ to $1.90$ --
physically impossible for European calls. The mechanism is
\emph{non-causality of the circular FFT}. The filter computes
\[
 \Delta\hat{W}_k
  = \bigl[\mathrm{IRFFT}\!\bigl(\tilde{H}\cdot\mathrm{RFFT}(z)\bigr)\bigr]_k,
\]
a circular convolution over the \emph{entire} horizon $N$. Consequently
$V_k = \xi_0 e^{\eta\hat{W}_k - \cdots}$ depends on future innovations
$z_{k+1},\ldots,z_{N-1}$ through the circular wrap-around, so $V_k$ and
$d\!\log S_{k+j}$ are spuriously correlated forward in time, violating
the martingale condition.

\paragraph{CausalSpectralFilter: architecture and initialisation.}
We replace the circular FFT with a causal lower-triangular Toeplitz
convolution:
\begin{equation}
 \Delta\hat{W}_k = \sum_{j=0}^{k} h_{k-j}\,z_j,
 \label{eq:causal-filter}
\end{equation}
where $h \in \mathbb{R}^N$ is a learnable real impulse response.
Because output at step $k$ depends only on $z_0,\ldots,z_k$, we have
$V_k \perp z_k$ and \eqref{eq:z-cholesky} is mathematically valid.
We implement \eqref{eq:causal-filter} as a single \texttt{F.conv1d}
with $(N-1)$ left-zero-padding, costing one matrix–vector product per
forward pass.

Initialisation uses the Riemann--Liouville increment kernel:
\begin{equation}
 h_0 = \frac{dt^H}{\Gamma(H+\tfrac{1}{2})},
 \qquad
 h_m = \frac{dt^H\bigl(g(m+1)-g(m)\bigr)}{\Gamma(H+\tfrac{1}{2})},
 \quad m \geq 1,
 \label{eq:rl-init}
\end{equation}
where $g(i)=i^{H-1/2}$. For $H < \tfrac{1}{2}$, $h_m < 0$ for $m \geq 1$
(roughness is encoded as negative memory in the innovations). We then
fine-tune with the same Bochner PSD loss used throughout, reaching a
final log-PSD MSE of $0.019$ (vs.\ $0.017$ for the circular filter) in
$3\,000$ epochs. The CausalSpectralFilter has $N = 252$ parameters
(vs.\ $N/2 = 126$ for the circular architecture) and an inference cost
of $12.6$~ms per $10\,000$ paths (vs.\ $1.9$~ms for the circular FFT)
on the same GTX 950 GPU, a $6.7\times$ slowdown that remains
$29\times$ faster than the Wood--Chan reference simulator. Hurst
recovery via structure function is preserved
($|\hat{H}-H| = 0.003$ for causal vs.\ $0.009$ for circular, both well
within the $\sqrt{\varepsilon}$ bound of \cref{cor:hurst}).

\paragraph{Results.}
Table~\ref{tab:causal-blp} and Figure~\ref{fig:causal-iv} summarise.
The causal filter with $z$-Cholesky achieves an IV RMSE of
\textbf{3.26~vp}, a \textbf{63$\times$ improvement} over the circular
architecture and better than the $\rho=0$ causal baseline (6.8~vp).
All ATM deltas are in $[0,1]$; the ATM bias shrinks from $-0.17$
($\rho=0$ baseline) to $-0.07$ ($z$-Cholesky), recovering $\sim\!60\%$
of the leverage-effect contribution.

\begin{table}[h]
\centering
\caption{IV RMSE vs BLP rBergomi reference ($\rho=-0.9$, 8k paths each)
and ATM $\Delta$ at $T=0.50$. Causal $z$-Cholesky is 63$\times$ better
than the circular architecture and reduces the delta bias by 60\% vs.\
the $\rho=0$ baseline.}
\label{tab:causal-blp}
\begin{tabular}{lrrrr}
\toprule
Condition & \multicolumn{3}{c}{IV RMSE (vp) per maturity} & ATM $\Delta$ \\
     & $T=0.25$ & $T=0.50$ & $T=1.00$ & $T=0.50$ \\
\midrule
BLP reference ($\rho=-0.9$)      &   &   &   & 0.689 \\
Circular SF  z-Cholesky ($\rho=-0.9$) & 57.7 & 89.8 & 359.6 & $>1$ \\
Causal SF   baseline  ($\rho=0$)  & 7.6 & 7.5 &  5.3 & 0.517 \\
Causal SF   z-Cholesky ($\rho=-0.9$) & \textbf{4.1} & \textbf{3.7} & \textbf{1.4} & \textbf{0.620} \\
\bottomrule
\end{tabular}
\end{table}

\begin{figure}[t]
 \centering
 \includegraphics[width=\textwidth]{causal_filter_comparison}
 \caption{%
  \textbf{CausalSpectralFilter vs circular SpectralFilter ($\rho=-0.9$).}
  Each panel shows one maturity.
  BLP reference (black) has the correct left skew.
  Causal $z$-Cholesky (green dashed) tracks the reference closely at
  all strikes; causal $\rho=0$ baseline (blue dotted) matches the IV
  \emph{level} but not the skew. Circular $z$-Cholesky (red
  dash-dot) is physically invalid (IVs exceed 100\% and deltas $>1$
  at long maturities, not shown on this scale).
 }
 \label{fig:causal-iv}
\end{figure}

\paragraph{Robustness across $H$.}
The improvement is not specific to $H = 0.1$.
Table~\ref{tab:causal-hurst} reports the sweep across the typical
rough-vol range $H \in \{0.05, 0.10, 0.15\}$. The causal $z$-Cholesky
IV RMSE stays in $[2.4, 3.5]$~vp across all $H$, while the circular
baseline exceeds $100$~vp at every $H$. The improvement factor
varies between $29\times$ and $60\times$, peaking at $H=0.1$ where the
circular failure is most severe. ATM delta bias decreases monotonically
with $H$ (from $-0.076$ at $H=0.05$ to $-0.050$ at $H=0.15$), consistent
with smoother $V$ paths being easier to match in distribution.

\begin{table}[h]
\centering
\caption{CausalSpectralFilter robustness sweep across $H$, all conditions
at $\rho=-0.9$, 8k paths. IV RMSE in vol points vs BLP reference;
$\Delta$ bias = $\hat\Delta_{\mathrm{ATM}} - \Delta_{\mathrm{ATM,ref}}$ at $T=0.50$.}
\label{tab:causal-hurst}
\begin{tabular}{ccccccc}
\toprule
$H$ & Circular SF & Causal SF & Causal SF & Improvement & ATM $\Delta$ bias \\
  & z-Cholesky & $\rho=0$ & z-Cholesky & factor   & (causal $z$-Chol) \\
\midrule
0.05 & 103.0 & 7.10 & 3.54 & $29\times$ & $-0.076$ \\
0.10 & 204.6 & 6.84 & \textbf{3.41} & $\mathbf{60\times}$ & $-0.073$ \\
0.15 & 106.1 & 6.83 & 2.39 & $44\times$ & $-0.050$ \\
\bottomrule
\end{tabular}
\end{table}

\paragraph{Closing the residual: matched-PSD training.}
The remaining $\sim\!-0.07$ delta bias of the causal $z$-Cholesky
surrogate is attributable to the V-distribution gap between the
RL-discretised BLP reference and the Wood--Chan-based fGn PSD used in
training. We test this directly by training a second causal filter on
the \emph{RL-fBm} PSD (estimated by Monte Carlo from $10^5$ paths of the
same RL discretisation that drives the reference). Table~\ref{tab:rl-match}
shows the result:

\begin{table}[h]
\centering
\caption{Effect of matching the training PSD to the reference V
distribution ($H=0.1$, $\rho=-0.9$, 5k surrogate paths). IV RMSE is
against the BLP reference; bias is
$\hat\Delta_{\mathrm{ATM}} - \Delta_{\mathrm{ATM,ref}}$. Matching the
training PSD to the RL-fBm distribution eliminates the residual bias.}
\label{tab:rl-match}
\begin{tabular}{lcccc}
\toprule
Training PSD & IV RMSE (vp) & Bias $T=0.25$ & Bias $T=0.50$ & Bias $T=1.00$ \\
\midrule
Wood--Chan (exact fBm)     & 3.42 & $-0.086$ & $-0.073$ & $-0.058$ \\
RL-fBm (matched to reference) & \textbf{2.57} & \textbf{$+0.011$} & \textbf{$+0.011$} & \textbf{$-0.007$} \\
\bottomrule
\end{tabular}
\end{table}

The matched-PSD surrogate achieves $2.57$~vp and an ATM delta bias of
$\lesssim 0.011$ at all maturities, essentially Monte Carlo noise.
This confirms the diagnosis: the leverage-channel residual was an
artefact of the training-vs-reference V-distribution mismatch, not a
deficiency of the causal architecture. The lesson generalises: for
any reference simulator, training the causal SpectralFilter on the
simulator's empirical PSD recovers $\rho \neq 0$ Greeks faithfully.

\paragraph{Cross-model: causal $z$-Cholesky on rough Heston.}
We finally test how much of the causal-architecture benefit transfers
to a \emph{cross-model} setting where the reference is rough Heston but
the surrogate uses rBergomi $V$. In §\ref{sec:greeks} this corresponds
to Surrogate~B (circular $z$-Cholesky), which failed at 188~vp with
$\Delta$ bias $+0.49$ (T=.50). Replacing the circular filter with the
causal architecture yields the headline numbers in Table~\ref{tab:rh-causal}.

\begin{table}[h]
\centering
\caption{Causal SpectralFilter applied to the rough Heston reference
($H=0.1$, $\kappa=2$, $\nu=1.9$, $\rho=-0.9$). Surrogate~$V$ paths are
rBergomi (not rough Heston); only the correlation architecture changes
from \cref{sec:greeks}'s Surrogate~B.}
\label{tab:rh-causal}
\begin{tabular}{lccc}
\toprule
Condition & IV RMSE (vp) & ATM $\Delta(T=0.50)$ & Bias \\
\midrule
Rough Heston reference            &     & 0.7745 &  \\
§\ref{sec:greeks} Surr~B (\emph{circular} z-Chol, $\rho=-0.9$) & 188  & ($>1$) & $+0.488$ \\
Bias~A floor (perfect-V, fresh BM)      &     & 0.5164 & $-0.258$ \\
Causal SF baseline ($\rho=0$)        & 8.0   & 0.5144 & $-0.260$ \\
Causal SF z-Cholesky ($\rho=-0.9$)      & \textbf{4.57} & \textbf{0.6247} & \textbf{$-0.150$} \\
\bottomrule
\end{tabular}
\end{table}

Two observations. First, the IV RMSE drops from $188$~vp
(§\ref{sec:greeks}'s circular Surr~B) to $4.57$~vp, a $41\times$
improvement, with all deltas physically in $[0,1]$, despite the
mismatched V dynamics. Second, and more striking, the causal
$z$-Cholesky bias of $-0.150$ is \emph{above} the $-0.258$ Bias~A
``floor'' that applies to surrogates with perfectly-matched $V$
marginals. This is not a contradiction: Bias~A measures the loss from
omitting the leverage channel entirely; the causal $z$-Cholesky
partially recovers it through the correlation architecture, recouping
$\sim\!42\%$ of the leverage contribution \emph{despite} using a
structurally different (log-normal rBergomi) $V$ distribution. The
remaining $-0.15$ is the irreducible mean-reversion gap that requires
joint $(V,S)$ training to close.

\medskip

\paragraph{Two negative results: spectral matching of rough Heston $V$
makes things worse.}
We attempted two natural variants of the causal SpectralFilter that
match rough Heston's $V$ marginal more closely:
\emph{(i)} train on the PSD of $V_{\mathrm{rH}}-\theta$ and use a linear
reconstruction $V_{\mathrm{surr}} = \theta + \mathrm{filter}(z)$;
\emph{(ii)} train on the PSD of $\log V_{\mathrm{rH}}$ and use the
rBergomi-style exp reconstruction
$V_{\mathrm{surr}} = \theta \exp(W - \tfrac{1}{2}\mathrm{Var}\,W)$.
Both \textbf{worsen} the cross-model bias (Table~\ref{tab:rh-negatives}).
Variant~(i) requires clipping $44\%$ of cells to enforce positivity,
inflating $\E[V_{\mathrm{surr}}]$ to $0.15$ ($3.7\times$ target).
Variant~(ii) collapses $\E[V_{\mathrm{surr}}]\to 0.004$ because the
empirical variance correction is dominated by
$\mathrm{Var}(d\log V_{\mathrm{rH}}) \approx 62$, log-$V$ excursions
when $V_{\mathrm{rH}}$ briefly visits very small values.

\begin{table}[h]
\centering
\caption{Attempts to match rough Heston $V$ within the spectral framework
all underperform the rBergomi-$V$ baseline. Spectral surrogates capture
second-order statistics; rough Heston $V$'s skewness and heavy tails
require higher-order modelling.}
\label{tab:rh-negatives}
\begin{tabular}{lccc}
\toprule
Surrogate $V$ design & IV RMSE (vp) & ATM $\Delta(T=.5)$ & Bias \\
\midrule
Causal SF, rBergomi $V$ (Surr~D above)       & \textbf{4.57} & \textbf{0.6247} & \textbf{$-0.150$} \\
Causal SF, linear $V = \theta + \mathrm{filter}(z)$ & 20.4 & 0.456 & $-0.319$ \\
Causal SF, $\log V$ matched, exp recon       & 9.7 & 0.485 & $-0.290$ \\
\bottomrule
\end{tabular}
\end{table}

Both failures reinforce the diagnosis: rough Heston's $V$ marginal is
too skewed and heavy-tailed for Gaussian second-order matching, and the
rBergomi log-normal $V$, \emph{despite a wrong distribution}, is the
best base process within the spectral family. Closing the remaining
$-0.15$ bias requires a Volterra-aware (non-spectral) architecture
that captures rough Heston's mean reversion and tail structure.

\paragraph{Volterra-aware prototype closes the cross-model bias.}
We prototype the prescribed architecture directly. The
\emph{VolterraSurrogate} is a learnable autoregressive module with
$N+3 = 255$ parameters, V-conditional diffusion, and a learnable RL
kernel:
\begin{equation}
 V_{k+1}
  = V_0 + c_d \!\!\sum_{j=0}^{k}\! w_{k-j}(\theta - V_j)
     + c_s \!\!\sum_{j=0}^{k}\! w_{k-j}\sqrt{V_j}\,z_j,
 \label{eq:volterra-surrogate}
\end{equation}
where $(c_d, c_s, \theta, \{w_i\}_{i=0}^{N-1})$ are trainable and
initialised at the rough Heston values
$c_d = \kappa\,dt^\alpha/\Gamma(\alpha)$,
$c_s = \nu\,dt^{\alpha-1/2}/\Gamma(\alpha)$,
$w_i = (i+1)^{\alpha-1}$, with $\alpha = H + \tfrac{1}{2}$.
Paired with the same $z$-Cholesky as before, this surrogate
\textbf{closes the cross-model bias to Monte Carlo noise}:

\begin{table}[h]
\centering
\caption{Volterra-aware prototype (\ref{eq:volterra-surrogate}) vs the
best spectral approach. Both use the same $z$-Cholesky correlation
($\rho=-0.9$). The V-conditional diffusion ($\sqrt{V_j}$ in the noise
term) is the architectural element that closes the cross-model bias.}
\label{tab:volterra-prototype}
\begin{tabular}{lccc}
\toprule
Surrogate & IV RMSE (vp) & ATM $\Delta(T=.5)$ & Bias \\
\midrule
Best spectral (causal SF, rBergomi $V$)  & 4.57 & 0.6247 & $-0.150$ \\
Volterra-aware ($\rho=0$)         & 8.99 & 0.5210 & $-0.254$ \\
\textbf{Volterra-aware ($z$-Chol, $\rho=-0.9$)} & \textbf{0.80} & \textbf{0.7765} & \textbf{$+0.002$} \\
\bottomrule
\end{tabular}
\end{table}

ATM delta bias at every maturity is $\lvert\hat\Delta - \Delta_{\mathrm{ref}}\rvert
\leq 0.005$, at the Monte Carlo noise floor with 5k surrogate paths.
The architecture validates the §\ref{sec:greeks}--§\ref{sec:causal}
diagnosis: the precise architectural element required is the
$\sqrt{V_j}$ coupling in the diffusion, which makes the update
\emph{nonlinear} in the past trajectory and cannot be expressed as a
spectral filter of $z$ alone.

\medskip

The Volterra-aware surrogate as prototyped is initialised at rough
Heston exactly; the result above is architectural. We also train the
same module end-to-end on observed data, starting from a heavily
misspecified initial state $(\theta, \kappa, \nu) = (0.10, 1.0, 0.5)$
versus targets $(0.04, 2.0, 1.9)$. Gradient flow through the
autoregressive Volterra update is stable with gradient clipping and
$V_{\min} = 10^{-6}$ to avoid $\sqrt{0}$ blow-ups; $200$ Adam steps
with batch $B = 128$.

Two loss formulations. \emph{Price-only}: $L = \|C_{\mathrm{surr}}
- C_{\mathrm{ref}}\|^2$ over the $(K, T)$ grid. This is feasibly
trainable but bottoms out at $5.81$~vp IV RMSE and $-0.088$ ATM bias
(Table~\ref{tab:volterra-trained}); call prices alone do not pin down
the $V$ dynamics tightly enough, many parameter triples give nearly
identical European call prices but different conditional behaviour.
\emph{Price + moment}: $L = \|C_{\mathrm{surr}} - C_{\mathrm{ref}}\|^2
+ \lambda \bigl( \|\mathbb{E}V_{\mathrm{surr}} - \mathbb{E}V_{\mathrm{ref}}\|^2
+ \|\mathrm{Var}\,V_{\mathrm{surr}} - \mathrm{Var}\,V_{\mathrm{ref}}\|^2 \bigr)$
with $\lambda = 50$. The moment terms are differentiable through the
same MC batch. This closes about half the remaining gap, reaching
\textbf{4.16~vp} and \textbf{$-0.054$} ATM bias, beating the spectral
floor on \emph{both} metrics for the first time using a learned
(not initialised) surrogate.

\begin{table}[h]
\centering
\caption{End-to-end training of the Volterra surrogate from misspecified
init $(\theta, \kappa, \nu) = (0.10, 1.0, 0.5)$. The auxiliary V-moment
loss resolves price-only identifiability and produces the first
\emph{trained} surrogate to beat the spectral floor.}
\label{tab:volterra-trained}
\begin{tabular}{lccc}
\toprule
Method (rough Heston, $\rho=-0.9$) & IV RMSE (vp) & ATM $\Delta(T=.5)$ & Bias \\
\midrule
Spectral best (causal SF rBergomi-V)    & 4.57 & 0.6247 & $-0.150$ \\
Volterra, RH-initialised (no training)   & 0.80 & 0.7765 & $+0.002$ \\
Volterra, price-only training        & 5.81 & 0.6869 & $-0.088$ \\
\textbf{Volterra, price + V-moment training} & \textbf{4.16} & \textbf{0.7208} & \textbf{$-0.054$} \\
\bottomrule
\end{tabular}
\end{table}

\paragraph{Sub-quadratic Volterra updates.}
The Volterra recursion above is $O(N^2)$ per path because the kernel
sum at step $k$ uses all of $V_0, \ldots, V_k$. Approximating the RL
kernel by a sum of $M$ exponentials,
$w_i \approx \sum_{m=1}^{M} a_m\, r_m^{\,i}$ with $r_m = e^{-\lambda_m}$,
turns the sum into $M$ recursively updatable states:
\begin{equation}
 S_m^X(k) = X_k + r_m\, S_m^X(k-1),
 \qquad \text{cost: } O(M) \text{ per step.}
 \label{eq:markov-state}
\end{equation}
We fit $(a_m, \lambda_m)$ from the rough-Heston RL kernel by
geometrically spacing the $\lambda_m$ over $[1/N, 1]$ and solving for
$a_m$ by nonneg least squares. The resulting
\emph{MarkovianVolterraSurrogate} has the same V-conditional structure
\eqref{eq:volterra-surrogate} but $O(MN)$ per path.
Table~\ref{tab:markov-volterra} shows the sweep:

\begin{table}[h]
\centering
\caption{Sub-quadratic Volterra via $M$-mode RL kernel approximation
($H=0.1$, $N=252$, batch $5\,000$, GPU). $M = 5$ already saturates the
kernel-fit error ($\approx 4\%$ relative) and gives the best IV-RMSE
of the sweep, $7\times$ faster than the exact $O(N^2)$ baseline.}
\label{tab:markov-volterra}
\begin{tabular}{lcccc}
\toprule
Surrogate & Fwd time (ms) & Speedup & IV RMSE (vp) & ATM $\Delta$ bias \\
\midrule
$O(N^2)$ Volterra (exact)      & 280 & $1\times$  & 0.91 & $-0.004$ \\
Markovian, $M=3$           & 47 & $6.0\times$ & 0.76 & $-0.001$ \\
\textbf{Markovian, $M=5$}      & \textbf{39} & \textbf{$7.2\times$} & \textbf{0.73} & \textbf{$+0.001$} \\
Markovian, $M=7$           & 40 & $7.1\times$ & 0.88 & $-0.008$ \\
Markovian, $M=10$          & 40 & $7.1\times$ & 0.85 & $-0.009$ \\
\bottomrule
\end{tabular}
\end{table}

All Markovian variants match the exact baseline within MC noise; the
$7\times$ speedup is bounded below by Python loop overhead and would
grow as $O(N/M)$ for longer time horizons. The Greeks are robust to a
$4\%$ relative kernel-fit error, a strong empirical indicator that
sub-quadratic approximations are practically usable.

\paragraph{Non-affine generalisation and full joint training.}
The diffusion term $\sqrt{V_j}$ in \eqref{eq:volterra-surrogate} is a
Heston-specific choice ($\sigma(V) = \nu V^{1/2}$). The same
autoregressive form accepts $\sigma(V) = V^\beta$ with $\beta$
learnable, call this the \emph{PowerVolterraSurrogate}. For
$\beta = 1$ this is a rough-lognormal-vol process; for $\beta = 0$ a
rough additive-vol process; intermediate $\beta$ interpolates.
Training the full parameter set
$(\theta, c_{\rm drift}, c_{\rm stoch}, \beta, \{w_i\})$ with the
price-plus-moment loss from a heavily misspecified initial state
$(\theta, \kappa, \nu, \beta) = (0.10, 1.0, 0.5, 0.8)$ \emph{and} a
$20\%$-perturbed kernel reaches IV RMSE $1.85$~vp and ATM delta bias
$+0.014$, closer to the RH-initialised $0.80$~vp ceiling than to
the spectral-floor $4.57$~vp. Final learned $\beta = 0.60$ (target
$0.5$). Individual parameter values do not match rough Heston, but
the joint trajectory and observables do, classical price-only
identifiability that the moment loss only partially resolves.

\paragraph{Live SPY option chain calibration (multi-maturity).}
We calibrate the joint Volterra surrogate to real SPY mid prices pulled
from Yahoo Finance at three expirations simultaneously: $T = 32$, $88$,
and $179$ days, with $S_0 = \$745.64$. All parameters
$(\theta, c_{\mathrm{drift}}, c_{\mathrm{stoch}}, \beta, \{w_i\})$
are \emph{shared} across the three maturities, only one $V$ path is
generated, and prices at each $T_i$ read from the appropriate column of
$V$. Twenty-six liquid OTM/ATM/ITM strikes total span
$K/S_0 \in [0.85, 1.15]$; market bid/ask spread implies $\sim\!1.4$~vp
IV uncertainty. Initialised at
$\theta_0 = (\mathrm{IV}^{\mathrm{ATM, short}})^2$ with $\rho = -0.85$,
the surrogate converges in $500$ Adam steps ($3$~min, GTX~950) to:

\begin{table}[h]
\centering
\caption{Multi-maturity calibration to live SPY options
($S_0 = \$745.64$, $\rho = -0.85$, $26$ calibration points across three
expirations). Overall IV RMSE $3.97$~vp; the short-end fits within
$2\times$ bid/ask uncertainty, while the long-end shows a systematic
under-pricing of OTM calls characteristic of single-kernel rough-vol
calibrations (the surrogate's variance grows too slowly with $T$).}
\label{tab:spy-multi}
\begin{tabular}{lccc}
\toprule
Expiry & $T$ (yr) & Strikes & IV RMSE (vp) \\
\midrule
2026-06-26 ($+32$d) & 0.088 & 8 & \textbf{2.27} \\
2026-08-21 ($+88$d) & 0.241 & 9 & 3.79 \\
2026-11-20 ($+179$d) & 0.490 & 9 & 5.17 \\
\midrule
\multicolumn{3}{l}{\textbf{Overall}} & \textbf{3.97} \\
\bottomrule
\end{tabular}
\end{table}

Final learned parameters
$(\theta, c_{\mathrm{stoch}}, \beta) = (0.037, 0.72, 0.52)$ recover
near-Heston dynamics ($\beta \approx 0.5$); the $\theta$ value is
consistent with the long-run SPX VIX level. The term-structure
residual at the $179$-day expiry, a $\sim 5$~vp uniform downward bias
across all strikes, reflects a known limitation of \emph{stationary-
parameter} rough-vol calibrations: a single kernel cannot match the
SPX ATM IV term structure (rising from $15.4\%$ at $32$d to $20.8\%$ at
$179$d here) without time-varying forward variance. Closing this is
the multi-maturity forward-consistency direction listed below.

\paragraph{SPX Hurst index from ATM term structure.}
Before the structural calibration, we estimate the implied Hurst exponent
directly from the SPX ATM IV term structure at three maturities
($T \in \{0.088, 0.241, 0.490\}$~yr, ATM IVs $15.6\%$, $18.1\%$, $20.0\%$).
Under rBergomi, $\mathrm{ATM\text{-}IV}(T) \propto T^{H-1/2}$ (Fukasawa
approximation), so a log-log OLS fit gives
$\hat{H}_\mathrm{SPX} = 0.145 \pm 0.002$ (standard error across 500
bootstrap resamples of the three-point fit).
We then train a SpectralFilter at $H = 0.145$ and evaluate against a BLP
reference at the same parameters: $H_\mathrm{est} = 0.151$
(structure-function, $|\Delta H| = 0.007$), IV RMSE $= 1.43$~vp
(Figure~\ref{fig:spx-hurst}), and variance ratio
$\mathrm{Var}(V_\mathrm{surr})/\mathrm{Var}(V_\mathrm{ref}) = 0.51$
(some variance is absorbed in the Itô correction).
The $|\Delta H| = 0.007$ is within the $\sqrt{\varepsilon}$ bound of
Corollary~\ref{cor:hurst} and confirms that the Cramér PINN can be
fitted directly to the market-implied Hurst index with no additional
calibration.

\begin{figure}[h]
 \centering
 \includegraphics[width=\textwidth]{spx_hurst_analysis}
 \caption{%
  \textbf{SPX Hurst analysis.}
  Left: log-log ATM IV vs.\ maturity; slope gives
  $\hat{H}_\mathrm{SPX} = 0.145$.
  Right: IV surface, SpectralFilter (trained at $H=0.145$, $\eta=1.485$)
  vs.\ BLP reference ($T\in\{0.25,0.50,1.00\}$~yr,
  $K/S_0\in[0.80,1.20]$).  IV RMSE $= 1.43$~vp.
 }
 \label{fig:spx-hurst}
\end{figure}

\paragraph{Forward-variance parameterisation: rBergomi-style.}
Multi-maturity SPY calibration with stationary $\theta$ (above) leaves
a uniform $\sim\!5$~vp residual at the long expiry, driven by SPX's
rising ATM-IV term structure ($15.4\% \to 20.8\%$ from $32$d to $179$d)
that a scalar $\theta$ cannot match. The cleanest structural fix is
the rBergomi forward-variance parameterisation:
\begin{equation}
 V_t \;=\; \xi(t) \cdot \exp\!\Bigl(\eta\, \hat W^H_t
     \;-\; \tfrac{1}{2}\eta^2\,\mathrm{Var}(\hat W^H_t)\Bigr),
 \label{eq:fwd-var}
\end{equation}
where $\hat W^H_t = \mathrm{cumsum}(\mathrm{CausalSF}(z))_t$ and $\xi(t)$
is a piecewise-linear forward-variance curve parameterised by knots at
the calibration maturities (log-parameterisation enforces positivity).
By the lognormal martingale identity, $\mathbb{E}[V_t] = \xi(t)$
exactly, so the ATM-IV term structure is pinned to $\sqrt{\xi(t)}$ by
construction. V is positive without clamping, the Feller-violation
pathology that broke the time-varying $\theta(t)$ Volterra attempt
vanishes.

Calibrating on the same SPY chain (with an empirical-mean renormalisation
that enforces $\mathbb{E}[V_t] = \xi(t)$ exactly at the finite batch),
we reach IV RMSE \textbf{$5.25$~vp} overall. The per-maturity breakdown
clarifies the residual:
\begin{table}[h]
\centering
\caption{rBergomi forward-variance calibration on SPY (same chain as
Table~\ref{tab:spy-multi}, $\rho=-0.85$, $\eta = 1.485$). Martingale
property holds exactly after the empirical-mean correction; remaining
residual is moneyness-driven, not maturity-driven.}
\label{tab:spy-rbergomi}
\begin{tabular}{lcccc}
\toprule
Maturity & ATM IV err & Deep-ITM err ($K/S_0 = 0.85$) & Deep-OTM err & Per-mat RMSE \\
\midrule
32d & $-1.2$~vp & $-13.7$~vp & $-5.1$~vp & 6.13 vp \\
88d & $-2.0$~vp & $-8.5$~vp & $-1.4$~vp & 4.05 vp \\
179d & $-4.8$~vp & $-8.7$~vp & $-2.1$~vp & 5.46 vp \\
\bottomrule
\end{tabular}
\end{table}

ATM strikes fit to within $\sim\!2$~vp; the residual is dominated by
\emph{deep ITM} ($K/S_0 \leq 0.90$) where the surrogate's smile is too
flat, \textbf{this is a skew mismatch, not a term-structure or
positivity issue.} A $\$2$ price error at $K=\$635$ (where the option
is mostly intrinsic value) is $10\%$ of the time value and only $2\%$
of the total price, but corresponds to $\sim\!14$~vp of IV error
because deep-ITM vega is tiny. Adam-with-uniform-price-MSE cannot see
this in its gradients.

\paragraph{Remaining open directions for Paper~3:}
The structural calibration story is now closed: stationary-$\theta$
Volterra ($3.97$~vp, no term structure), time-varying-$\theta$
Volterra ($5.59$~vp, Feller violation), and rBergomi forward variance
($5.25$~vp, no Feller, exact martingale). Closing the remaining
moneyness-driven residual requires:
(i) IV-direct (vega-weighted or implicit-differentiation) loss to make
the optimiser see deep-strike error;
(ii) two-stage calibration (ATM curve first, then smile) so $\eta$ can
be pushed to SPX-realistic levels without breaking ATM fit;
(iii) richer auxiliary moment losses (autocovariance, third moments) to
constrain the joint $(V, S)$ law beyond second-order matching.

% ═════════════════════════════════════════════════════════════════════════════
\section{Discussion}
\label{sec:discussion}
% ═════════════════════════════════════════════════════════════════════════════

\paragraph{Why the exact PSD matters.}
The power-law PSD $S_H(f) \sim C|f|^{1-2H}$ is the correct
\emph{asymptotic} behaviour of fGn for $f \to 0$. At mid-frequencies
the exact PSD deviates from the power law; for $H = 0.1$ this deviation
exceeds 10\% relative error across most of the frequency range.
Matching the asymptotic shape correctly recovers the power-law slope but
mis-trains the correlation structure at short lags, biasing long-maturity
option prices (the $T = 1$ column in Table~\ref{tab:iv-rmse}).
The PSD phase diagram (Figure~\ref{fig:psd-phase}) quantifies the
regime boundary: exact PSD training is essential below $H\approx 0.15$;
above $H=0.20$ both targets converge to comparable IV RMSE. The
relative $L^2(\log)$ PSD distance drops from $0.072$ at $H=0.05$
to $0.011$ at $H=0.30$, matching the observed IV RMSE gap.
Horizon sensitivity (Table~\ref{tab:horizon-sweep}) shows this regime
boundary also applies in the time domain: longer observation windows
help Hurst recovery but amplify the low-frequency cutoff penalty,
making the Cramér PINN most reliable for $T \leq 2$~yr.

\paragraph{BochnerMLP: architecture-agnostic spectral loss.}
The MLP surrogate outperforms SpectralFilter on overall RMSE despite
generating non-Gaussian paths. This is expected: the MLP has more
capacity (41k vs.\ 126 parameters) and can approximate the fGn PSD more
closely. However, non-Gaussianity means the Itô correction
$-\frac{1}{2}\eta^2 t^{2H}$ does not exactly ensure $\E[\hat{V}_t] =
\xi_0$ (Jensen's inequality with non-Gaussian $\hat{W}_t$); we use the
empirical variance correction \eqref{eq:vhat} to mitigate this.
SpectralFilter, being a linear filter of Gaussian input, is
\emph{guaranteed} to produce Gaussian fGn paths and avoids this issue
entirely.

\paragraph{Application to deep hedging.}
\citet{buhler2019deep} require large batches of simulated price paths for
training a hedging neural network. The spectral surrogate reduces
path generation from $O(N^2)$ (Cholesky-based fBm) or $O(N \log N)$
(Wood--Chan) to a single forward pass of the trained network, enabling
on-the-fly batch sampling. The $\approx 1.2$~vp IV accuracy is well
within the calibration precision of typical options desks.
However, \cref{sec:greeks} shows that Greeks require the joint $(V, S)$
law: the spectral surrogate, which matches only the marginal distribution
of $V$, carries an irreducible delta bias of $\approx -0.27$ at ATM
when $\rho = -0.9$. Deep hedging with a spectral surrogate therefore
requires either a joint-law surrogate (future work) or a Malliavin
score correction to recover the leverage contribution. A circular-FFT $z$-Cholesky attempt fails (204~vp, delta $>1$) due to
wrap-around making $V_k$ depend on future innovations
(Section~\ref{sec:greeks}). The CausalSpectralFilter resolves this
exactly, achieving 3.26~vp and delta bias $-0.07$, a 63$\times$ RMSE
improvement available at no additional training cost.

\paragraph{Comparison to Markovian approximations.}
A standard industry approach to rough volatility is to approximate the
fractional kernel by a finite sum of exponentials, yielding a
\emph{Markovian} multi-factor model that can be simulated with
conventional Euler schemes. Two leading constructions are:

\begin{itemize}
 \item \textbf{Lifted Heston} \citep{abijaber2019affine,el2019perfect}:
  approximates the fractional Riccati--Volterra equation by an
  $n$-dimensional affine jump-diffusion, with $n=5$--$20$ factors
  needed for $H \approx 0.1$. Pricing is via Fourier inversion
  (characteristic function available in closed form), but the
  characteristic function becomes increasingly expensive as $n$ grows.
 \item \textbf{3-factor Bergomi} \citep{bayer2020markovian}:
  parameterises a low-dimensional Markovian system whose joint
  (stock, variance) dynamics approximate the rough Bergomi
  smile. Calibration requires non-convex optimisation over
  $\sim 10$ parameters and the approximation degrades for
  $H < 0.07$, where 3 factors cannot capture the roughness.
\end{itemize}

These Markovian methods share a fundamental limitation: the number of
factors controls the quality of the fractional-kernel approximation, and
the approximation error propagates to option prices in a way that depends
on both the strike and the maturity. In contrast, the spectral surrogate
\emph{directly targets the PSD} of the fGn increments, bypassing the
kernel-approximation step entirely. The implication is that the spectral
surrogate's accuracy is determined by the spectral loss convergence
(Theorem~\ref{thm:main}) rather than by a factor-count parameter, and it
does not degrade as $H \to 0$, a regime where Markovian approximations
require increasingly many factors or break down altogether (Table~\ref{tab:hurst}
shows $|\Delta H| \leq 0.013$ at $H=0.05$, where $n=3$ Markovian
approximations typically incur bias $>0.03$ in recovered $H$).

A secondary advantage is that the spectral surrogate produces paths in
a single forward pass with no matrix exponentials, no Riccati ODE
integration, and no Fourier inversion, it is a pure feed-forward
generator, making it trivially vectorisable and GPU-friendly. The
trade-off is a one-time training cost and the $W_2$ convergence
guarantee that holds only for the Gaussian projection of the law
(Theorem~\ref{thm:main}, Assumption~\ref{ass:gauss}).

\paragraph{Connection to Paper 1.}
The proof structure is a one-assumption extension of \citet{TODO-paper1}:
Bochner (stationary levels) $\to$ Cramér (stationary increments);
all other lemmas carry over unchanged (see the proof-extension guide,
\texttt{proof\_extension\_guide.tex}). The new content is the Hurst
corollary (\cref{cor:hurst}) and the financial validation.

% ═════════════════════════════════════════════════════════════════════════════
\section{Conclusion}
\label{sec:conclusion}
% ═════════════════════════════════════════════════════════════════════════════

We have introduced the Cramér PINN, a neural-network surrogate for the
rough Bergomi and rough Heston models trained to match the exact power
spectral density of fractional Gaussian noise. The main contributions
are: a $W_2$ convergence theorem extending the Bochner PINN framework to
processes with stationary increments; a Hurst exponent consistency
corollary; an empirical demonstration that exact PSD matching achieves
1.21--1.33 vol-point IV surface RMSE under $\rho = 0$, $1.7\times$
better than the asymptotic power-law approximation and near the Monte
Carlo noise floor; a Greek bias analysis showing that the marginal $V$ surrogate carries
an irreducible ATM delta bias of $-0.27$ ($35\%$) when $\rho = -0.9$
under \emph{rough Heston}; a \emph{CausalSpectralFilter} (causal
convolution, same PSD loss) that resolves the circular-FFT wrap-around
pathology, achieving 3.26~vp vs.\ the BLP rBergomi reference (or
\textbf{2.57~vp} with matched-PSD training, leaving only Monte-Carlo
noise) and reducing the rBergomi delta bias to $\lesssim 0.01$; on rough
Heston the same spectral architecture cuts the cross-model bias from
$+0.49$ to $-0.15$ (41$\times$ IV RMSE improvement), and a
Volterra-aware autoregressive surrogate closes the residual to MC noise
(0.80~vp, bias $\leq 0.005$); a Hurst estimator comparison showing that
structure function achieves the smallest bias ($|\text{bias}|\leq 0.011$)
for $H\in[0.05,0.30]$ while PSD slope and R/S are unreliable below $H=0.2$;
a PSD phase diagram establishing $H\approx 0.15$ as the threshold below
which exact (not asymptotic) PSD training is essential; an observation-horizon
sweep recommending $T\leq 2$~yr for IV RMSE below $1.6$~vp; and a live SPX
Hurst estimate $\hat{H}_\mathrm{SPX}=0.145$ from the ATM IV term structure,
recovered by the surrogate to within $|\Delta H|=0.007$.

The CausalSpectralFilter introduced in Section~\ref{sec:greeks} already
handles $\rho \neq 0$ for rBergomi (3.26~vp vs.\ BLP reference).
Remaining natural extensions are: (i) a joint $(V, S)$ surrogate trained
on paired paths, which would close the residual $-0.07$ delta bias
attributable to the V-distribution gap; (ii) the deep hedging application
\citep{buhler2019deep}, for which the causal surrogate is a drop-in path
generator; and (iii) extension to other rough volatility models (RFSV,
multifactor Bergomi) where the same Cramér PSD loss applies.

% ═════════════════════════════════════════════════════════════════════════════
\appendix
\section{Proofs}
% ═════════════════════════════════════════════════════════════════════════════

\subsection{Proof of Theorem~\ref{thm:main}}

We bound the three terms in \eqref{eq:main-bound} separately and then
combine. Throughout, Assumptions~\ref{ass:si}--\ref{ass:gauss} are in force.

\paragraph{Reduction to the Gelbrich formula.}
Under Assumption~\ref{ass:gauss}, both $u_\theta$ and $u_{\mathrm{ref}}$
are centred Gaussian processes; their laws are completely determined by
their covariance operators $C_\theta, C_{\mathrm{ref}}$ on $L^2([0,T])$.
The Bures--Gelbrich formula \citep{gelbrich1990formula} gives
\begin{equation}\label{eq:gelbrich-main}
 W_2^2\!\bigl(\mathcal{L}(u_\theta),\,\mathcal{L}(u_{\mathrm{ref}})\bigr)
 = \bigl\|C_\theta^{1/2} - C_{\mathrm{ref}}^{1/2}\bigr\|_{\HS}^2.
\end{equation}
We bound the right-hand side in four steps.

\paragraph{Step 1 (Square-root Lipschitz).}
Because both processes have stationary increments with Cram\'er spectral
densities satisfying Assumption~\ref{ass:nd}, the covariance operators are
bounded below: $C_\theta, C_{\mathrm{ref}} \geq \mu_{\min} I$ on the
spectral support $[\omega_0,\Omega]$. The operator square-root Lipschitz
bound \citep[Thm.~X.1.1]{bhatia1997matrix} then gives
\begin{equation}\label{eq:sqrt-lip}
 \bigl\|C_\theta^{1/2} - C_{\mathrm{ref}}^{1/2}\bigr\|_{\HS}
 \;\leq\;
 \frac{1}{2\sqrt{\mu_{\min}}}\,
 \bigl\|C_\theta - C_{\mathrm{ref}}\bigr\|_{\HS},
\end{equation}
so that $W_2^2 \leq (4\mu_{\min})^{-1}\|C_\theta - C_{\mathrm{ref}}\|_{\HS}^2$.

\paragraph{Step 2 (Cram\'er kernel to Hilbert--Schmidt norm).}
By Assumption~\ref{ass:si}, the covariance kernel difference is
\[
 \Delta K(t,s) = \int_{\omega_0}^{\Omega}
  (e^{i\omega t}-1)\,\overline{(e^{i\omega s}-1)}\,
  \bigl(S_\theta(\omega)-S_{\mathrm{ref}}(\omega)\bigr)\,d\omega.
\]
Applying Fubini's theorem, $|e^{i\omega t}-1|^2 \leq 4$, and
Cauchy--Schwarz on $[\omega_0,\Omega]$:
\begin{equation}\label{eq:parseval-main}
 \bigl\|C_\theta - C_{\mathrm{ref}}\bigr\|_{\HS}^2
 = \int_0^T\!\!\int_0^T |\Delta K(t,s)|^2\,dt\,ds
 \;\leq\;
 \frac{T^2}{2\pi}\,
 \bigl\|S_\theta - S_{\mathrm{ref}}\bigr\|_{L^2([\omega_0,\Omega])}^2.
\end{equation}
(The factor $2\pi$ absorbs the exact integral
$\int_0^1|e^{i\omega t}-1|^2\,dt = 2-2\sin\omega/\omega \leq 4$.)

Combining Steps~1--2 with \eqref{eq:gelbrich-main}:
\begin{equation}\label{eq:spectral-bound}
 W_2\!\bigl(\mathcal{L}(u_\theta),\,\mathcal{L}(u_{\mathrm{ref}})\bigr)
 \;\leq\;
 \frac{T}{2\sqrt{2\pi\mu_{\min}}}\,
 \bigl\|S_\theta - S_{\mathrm{ref}}\bigr\|_{L^2([\omega_0,\Omega])}
 \;=:\;
 \frac{C_T}{\mu_{\min}^{1/2}}\,
 \bigl\|S_\theta - S_{\mathrm{ref}}\bigr\|_{L^2([\omega_0,\Omega])}.
\end{equation}
This is the spectral term in \eqref{eq:main-bound} with $C_T = T/(2\sqrt{2\pi})$.

\paragraph{Step 3 (Discretisation error).}
The spectral loss $\mathcal{L}(\theta)$ is evaluated on a discrete grid of
$N_s$ frequency points in $[\omega_0,\Omega]$. The surrogate fGn PSD
$S_\theta$ is a piecewise-smooth function inheriting the Sobolev
regularity $H^{H-\varepsilon}$ of fBm increments ($\beta_H =
\tfrac{1}{2}\min(H-\tfrac{1}{2}+\tfrac{d}{2},1)$ for dimension $d$).
Standard FFT quadrature theory gives
\[
 \bigl\|S_\theta^{\mathrm{discrete}} - S_\theta^{\mathrm{continuous}}
 \bigr\|_{L^2([\omega_0,\Omega])}
 \;\leq\; C_{\mathrm{quad}}\,N_s^{-\beta_H},
\]
with $C_{\mathrm{quad}}$ depending on $\|S_\theta\|_{H^{\beta_H}}$
and the geometry of the frequency grid
(cf.\ \citet{TODO-paper1}, Appendix~A, Lemma~A.4).
Inserting into \eqref{eq:spectral-bound} and setting
$C_2 = C_T\,C_{\mathrm{quad}}/\mu_{\min}^{1/2}$ gives the
discretisation term $C_2\,N_s^{-\beta_H}$.

\paragraph{Step 4 (Low-frequency cutoff).}
The spectral window $[\omega_0,\Omega]$ excludes the interval $[0,\omega_0)$
where the fBm PSD $S_H(\omega) = C_H|\omega|^{1-2H}$ diverges for $H < 1/2$.
Restricting to $\omega_0 > 0$ is necessary for Assumption~\ref{ass:nd} to
hold. The penalty for this restriction enters via $\mu_{\min}$: since
\[
 \mu_{\min} \;=\; S_H(\omega_0) \;=\; C_H\,\omega_0^{1-2H},
\]
the constant in \eqref{eq:spectral-bound} satisfies
$\mu_{\min}^{-1/2} = C_H^{-1/2}\,\omega_0^{-(1-2H)/2}$.
For any fixed spectral gap $\|S_\theta-S_{\mathrm{ref}}\|_{L^2}$, the
$W_2$ bound degrades as $\omega_0^{-(1-2H)/2}$ as the window shrinks toward
the origin. Setting
\[
 C_3\,\omega_0^{-(1-2H)/2}
 := \frac{C_T\,\|S_\theta-S_{\mathrm{ref}}\|_{L^2}}{C_H^{1/2}}\,
   \omega_0^{-(1-2H)/2}
\]
identifies the third term; it vanishes as $\omega_0 \to 0$ only when
$H > 1/2$ (the PSD then integrates to zero at the origin).

\paragraph{Combining the three terms.}
The triangle inequality on $W_2$ applied through the spectral estimate at
grid resolution $N_s$ and at frequency cutoff $\omega_0$ gives
\begin{align*}
 W_2\!\bigl(\mathcal{L}(u_\theta),\mathcal{L}(u_{\mathrm{ref}})\bigr)
 &\leq\;
 \underbrace{\frac{C_T}{\mu_{\min}^{1/2}}
 \bigl\|S_\theta-S_{\mathrm{ref}}\bigr\|_{L^2([\omega_0,\Omega])}}_{\text{spectral term}}
 \;+\;
 \underbrace{C_2\,N_s^{-\beta_H}}_{\text{discretisation}}
 \;+\;
 \underbrace{C_3\,\omega_0^{-(1-2H)/2}}_{\text{cutoff}}.
\end{align*}
This is \eqref{eq:main-bound}. The three contributions are independent:
the spectral term vanishes when the trained PSD matches the target;
the discretisation term vanishes as the frequency grid is refined;
the cutoff term reflects the irreducible cost of avoiding the $\omega = 0$
singularity.\hfill$\square$

\subsection{Wood--Chan simulation}
\label{app:wood-chan}

Exact simulation of fGn using circulant embedding
\citep{wood1994simulation, dietrich1997fast}.

The method exploits the fact that the covariance matrix of fGn increments
$(\Delta W^H_1, \ldots, \Delta W^H_N)$ is a Toeplitz matrix
$\Gamma_{jk} = \tfrac{1}{2}(|j-k+1|^{2H} - 2|j-k|^{2H} + |j-k-1|^{2H})$.
Every Toeplitz matrix can be embedded in a circulant $C$ of size $2N$;
if the embedding is positive semidefinite, the square root
$C^{1/2}$ is available via FFT in $O(N\log N)$. A realisation
of fGn is then $\Delta W^H = \mathrm{Re}(\mathcal{F}^{-1}[C^{1/2}z])$
where $z \sim \mathcal{CN}(0, I)$. Davies and Harte (1987) showed the
embedding is always positive semidefinite for $H \in (0,1)$ after
mirroring; \citet{wood1994simulation} and \citet{dietrich1997fast}
provide the practical algorithm. Our implementation is in
\texttt{rbm\_reference.py}, function \texttt{simulate\_fBm\_woodchan}.

% ─────────────────────────────────────────────────────────────────────────────
\bibliographystyle{plainnat}
\bibliography{references}
\end{document}
